On prélève un volume $V_{1}=\qty{50}{\mL}$ d'une solution d'acide chlorhydrique
de concentration $c_{1}=\qty{0.50}{\mol\per\L}$, ainsi qu'un volume
$V_{2}=\qty{20}{\mL}$ d'une seconde solution d'acide chlorhydrique où la
concentration des ions oxonium vaut~:
$\left[\ce{H3O+}\right]_{2}=\qty{0.20}{\mol\per\L}.$

On mélange les deux volumes prélevés.
\begin{enumerate}
    \item Calculer les quantités de matière d'ions oxonium présents dans chaque
          prélèvement~; en déduire la concentration $\left[\ce{H3O+}\right]$
          dans le mélange.
    \item Procéder aux mêmes calculs que précédemment pour les ions chlorure.
    \item On remplace le volume $V_{2}$ par un volume $V_{3}=\qty{50}{\mL}$
          d'une solution de chlorure de sodium de concentration
          $c_{3}=\qty{0.25}{\mol\per\L}$.
    
          Calculer les concentrations de tous les ions présents dans le mélange
          obtenu.
\end{enumerate}

